\section{Introduction}\label{sec:introduction}
A service agreement can return to the same public operating condition after different sequences of renewal reviews. Those histories need not permit the same future service. A customer who could have exited at an earlier review may have accepted a different remaining payment obligation. A provider that remembers only the current public condition can therefore lose value even when demand is Markovian, prices are fixed, and neither party has private information. The operational question is not whether full history can improve one schedule. It is whether the accepted optimum admits a compact exact representation, how large that representation must be in the worst case, and how much additional writable information exact execution requires.

We answer these questions for positive additive payments, separable strictly concave quadratic operating rewards, interval controls, and an exogenous recombining public graph. The exact compiler excludes intertemporal switching costs and coupled vector commitments. These conditions make one accumulated payment price sufficient for optimal continuation. Each participation cap reflects that price at a vertex-specific barrier. We construct the complete parameterized response family before choosing a root promise; the construction neither obtains an extensive-form optimum first nor enumerates a scenario tree.

The classical starting point is explicit. After probability weighting, the unfolded problem is separable concave resource allocation with laminar descendant constraints, in the lineage of \citet{Mjelde1983} and \citet{Tang1990}. Normalized memoization of that occurrence-indexed reduction produces the same public-graph recurrence. Parametric solution paths and sensitivity are also established themes: \citet{KlimmWarode2021} give output-polynomial algorithms for piecewise-quadratic parametric flow, while \citet{HarksKlimmPeis2018} study sensitivity of separable optimization over integral polymatroids. We therefore claim neither scalar marginal allocation, memoization, generic piecewise-affine parametricity, nor state minimization as new principles. The residual contribution is quantitative and representation-specific: on the compact public graph, the complete accepted-service response has at most $3N$ distinct knots, this linear knot bound and the resulting quadratic response-storage order are both worst-case tight, and the coefficient-event implementation has polynomial rational bit complexity.

The second contribution concerns execution rather than compilation. The reached price graph is a deterministic output transducer driven by realized public transitions. Backward state minimization is classical for deterministic machines and related state-aggregation models \citep{BeanBirgeSmith1987,GivanDeanGreig2003,Peyriere2023}. Applying that principle to the reached transducer yields the minimum equivalent behavior partition. The contract-specific statements are that a continuous-control problem generates a finite reached machine; its behavioral classes are contiguous in price order; the additional writable alphabet equals the largest within-vertex class count; distinct numerical prices can collapse to one behavior; and a renewal family attains the linear worst-case memory order. We report separately the additional writable alphabet, the number of reached public-state/class pairs, and the read-only response representation and coefficient precision.

The renewal family also gives an exact value--memory frontier with heterogeneous branch probabilities and operating costs. Along every strictly ordered mean-preserving radial spread of continuation caps, insufficient deterministic memory has strictly increasing loss, while the exact alphabet requirement can jump at arbitrarily small spread. The statement is a radial comparative static, not a claim of monotonicity under every convex-order or majorization relation. The randomized exact-optimum lower bound uses an explicit information-accounting convention; pathwise feasibility with respect to private controller randomization only strengthens that lower bound.

The service interpretation is positioned against finite policy graphs in dynamic contracting. \citet{Zhang2012} uses finite policy graphs to approximate continuation contracts in an infinite-horizon adverse-selection model. Our model has public information and a finite horizon, and it does not solve adverse selection. Its policy graph is instead an exact compiler-generated transducer for a fixed root query. This distinction, together with the public-graph promise-grid comparator, prevents a graph representation itself from carrying the novelty claim.

The computational study is organized around the residual theorem claims. An independently implemented occurrence-indexed laminar reduction checks exact equality; a public-graph promise method checks a competing compact state representation; dense response families expose output size; and repeated root inversions measure the use of a compiled parameterized path. Timing against a deliberately nonmemoized tree is retained only as an implementation comparison, not as evidence against the equivalent normalized memoized reduction or against specialized parametric algorithms.

The broader accepted-adaptation results on switching, comparator-adjusted rents, continuous states, and certified deployment remain available in the retained theory volumes with their statements and proofs unchanged. They are not used to enlarge the domain of the finite-price theorem. The paper's contribution is therefore a sharp exact theorem package for the stated subclass: a compact quotient, tight response complexity, explicit bit complexity, and a precisely accounted minimal writable implementation.